\documentclass[aspectratio=169]{beamer}
\usetheme{Madrid}
\usecolortheme{default}
\setbeamertemplate{navigation symbols}{}
\setbeamertemplate{footline}[frame number]

\usepackage[T1]{fontenc}
\usepackage[utf8]{inputenc}
\usepackage{lmodern}
\usepackage{amsmath,amssymb,mathtools,bm}
\usepackage{booktabs}
\usepackage{array}

\title{Reinforcement Learning for Quantum Control}
\author{Morten Hjorth-Jensen}
\date{March 17, 2026}

\begin{document}

\begin{frame}
  \titlepage
\end{frame}

\begin{frame}{Motivation}
\begin{itemize}
  \item Quantum sensors permit fine-grained control over probe evolution and measurement settings.
  \item The central design problem is sequential: which control should be applied \emph{next}, given what has already been learned?
  \item For realistic platforms, the estimation stack contains many non-smooth components: sampling of measurement outcomes, particle-filter resampling, and optional weak-measurement backaction.
  \item The article proposes a \alert{model-aware RL} framework that treats experimental design as a trainable control problem while retaining a faithful physical model.
\end{itemize}
\vspace{0.4em}
\textbf{Talk focus:} the technical machinery in Appendices A--I, with figures intentionally omitted so they can be inserted later.
\end{frame}

\begin{frame}{Main thesis of the paper}
\begin{block}{Core idea}
Use a differentiable simulation of the experiment, Bayesian inference by particle filtering, and policy optimization to learn adaptive or non-adaptive control laws for quantum metrology.
\end{block}
\begin{itemize}
  \item \textbf{Bayesian mode:} optimize the final estimation error after a complete measurement sequence.
  \item \textbf{Frequentist mode:} optimize the Cramer-Rao bound through the Fisher information matrix.
  \item \textbf{Model-aware:} gradients are propagated through a known sensor model rather than inferred only from trial-and-error interaction.
  \item \textbf{Applications in the paper:} NV-center DC magnetometry and the agnostic Dolinar receiver.
\end{itemize}
\end{frame}

\begin{frame}{The measurement loop as a control architecture}
\begin{enumerate}
  \item From the current posterior summary, the agent outputs the next control $x_{t+1}$.
  \item The probe evolves and is measured, producing a stochastic outcome $y_{t+1}$.
  \item The Bayesian posterior is updated and fed back to the agent.
\end{enumerate}
\vspace{0.5em}
Adaptive policies are written as
\[
  x_{t+1}=F_{\lambda}\!\big(P(\theta\mid x_t,y_t);\, y_t;\, R_t;\, t\big),
\]
while non-adaptive strategies suppress dependence on the posterior and outcomes:
\[
  x_{t+1}=F_{\lambda}(R_t,t).
\]
\begin{itemize}
  \item $\lambda$: trainable parameters of the policy (typically a neural network).
  \item $R_t$: resources consumed up to step $t$ (time, photons, signal amplitude, \ldots).
\end{itemize}
\end{frame}

\begin{frame}{Appendix A: probe, controls, and measurement model}
A generic experiment is cast as
\[
  \rho \longmapsto \rho_{x,\theta}=\mathcal E_{x,\theta}(\rho),
\]
where $\rho$ is a reference probe state, $x$ denotes tunable controls, and $\theta\in\Theta$ are the unknown parameters.
\vspace{0.6em}
\begin{itemize}
  \item Controls can be continuous or discrete, and may represent durations, detunings, phases, or measurement settings.
  \item Parameters of interest can coexist with \emph{nuisance parameters}; estimating the latter can still improve the former.
  \item A measurement is modeled by a POVM $\mathcal M^x=\{M^x_y\}$ with likelihood
  \[
    p(y\mid \theta,x)=\mathrm{tr}\!\left(M^x_y\rho_{x,\theta}\right).
  \]
\end{itemize}
\end{frame}

\begin{frame}{Appendix A: projective versus weak measurements}
\begin{itemize}
  \item For \textbf{projective measurements}, the probe is reinitialized after each shot; the next likelihood depends only on $x_{t+1}$ and $\theta$.
  \item For \textbf{weak measurements}, information remains in the probe state, so the likelihood depends on the whole trajectory:
  \[
    p(y_{t+1}\mid x_{t+1},y_t,\theta)=\mathrm{tr}\!\left(M^{x_{t+1}}_{y_{t+1}}\,\rho_{x_t,y_t,\theta}\right).
  \]
  \item This turns the estimator into a partially observed control problem in which the posterior and probe state both carry memory.
\end{itemize}
\vspace{0.6em}
\textbf{Control lesson:} the state fed to the agent is not the raw experimental history, but a compressed belief-state summary derived from Bayesian inference.
\end{frame}

\begin{frame}{Appendix B: particle-filter representation of the posterior}
The posterior is approximated by a weighted ensemble
\[
  P(\theta)\approx\sum_{j=1}^N w_j^t\,\delta(\theta-\theta_j),
\]
with particles $\{\theta_j\}_{j=1}^N$ and weights $\{w_j^t\}_{j=1}^N$.
\vspace{0.4em}
\begin{itemize}
  \item Initialize particles from the prior $\pi(\theta)$ and set $w_j^0=1/N$.
  \item Apply Bayes updates after each observation:
  \[
    w_j^{t+1}=\frac{p(y_{t+1}\mid x_{t+1},\theta_j)w_j^t}{\sum_{i=1}^N p(y_{t+1}\mid x_{t+1},\theta_i)w_i^t}.
  \]
  \item The posterior mean and covariance become
  \[
    \hat\theta_t\approx\sum_j w_j^t\theta_j,
    \qquad
    \Sigma_t\approx\sum_j w_j^t(\theta_j-\hat\theta_t)(\theta_j-\hat\theta_t)^\top.
  \]
\end{itemize}
These moments form the default summary presented to the policy.
\end{frame}

\begin{frame}{Appendix B.3: resampling, soft resampling, and proposal moves}
\small
When weights collapse onto a few particles, the filter is refreshed by resampling.
\vspace{0.25em}
\begin{block}{Soft resampling}
Instead of sampling directly from $w_j$, define
\[
  q_j=\alpha w_j + (1-\alpha)\frac{1}{N},
\]
and draw new particles $\theta'_j=\theta_{\phi(j)}$ from $q_j$.
\end{block}
Corrected weights preserve the represented posterior:
\[
  w'_j\propto \frac{w_{\phi(j)}}{q_{\phi(j)}}.
\]
Further steps in Appendix B:
\begin{itemize}
  \item \textbf{Gaussian perturbation:} move particles toward $\hat\theta_t$ and break degeneracy.
  \item \textbf{New-particle proposal:} inject particles from $\mathcal N(\hat\theta_t,\Sigma_t)$ to raise local resolution.
\end{itemize}
\normalsize
\end{frame}

\begin{frame}{Appendix B.4: state particle filter for weak measurements}
For weak measurements, each parameter particle carries its own hypothetical probe state:
\[
  \rho_{\theta_j,\tau_t}=\mathcal M^{x_t}_{y_t}\circ\mathcal E_{\theta_j,x_t}\circ\cdots\circ
  \mathcal M^{x_1}_{y_1}\circ\mathcal E_{\theta_j,x_1}(\rho).
\]
The normalized measurement backaction map is
\[
  \mathcal M^x_y(\rho)=\frac{M^x_y\rho M^{x\dagger}_y}{\mathrm{tr}[M^x_y\rho M^{x\dagger}_y]}.
\]
\begin{itemize}
  \item The state particle filter stores $(\theta_j,w_j,\rho_{\theta_j,\tau_t})$.
  \item It enables trajectory-dependent likelihood evaluation for probes that are not reinitialized between measurements.
  \item This is the natural extension from Bayesian filtering over parameters to Bayesian filtering over \emph{parameter-conditioned states}.
\end{itemize}
\end{frame}

\begin{frame}{Why differentiation is hard}
The training objective depends on operations that are stochastic or discrete:
\begin{itemize}
  \item measurement outcomes $y_t\sim p(y_t\mid x_t,\theta)$,
  \item categorical particle resampling,
  \item stochastic perturbation and proposal draws,
  \item optional discrete decision outputs in classification tasks.
\end{itemize}
\vspace{0.5em}
Automatic differentiation treats sampled values as constants in the backward pass, so naively
\[
  \frac{d y_t}{d\lambda}=0,
\]
and the gradient would fail to capture how the policy changes the distribution of future trajectories.
\vspace{0.5em}
\textbf{Appendices C and D} supply the fix: combine score-function terms with differentiable surrogates wherever possible.
\end{frame}

\begin{frame}{Appendix C: differentiable particle filtering}
Three ingredients make the particle filter trainable.
\vspace{0.3em}
\begin{enumerate}
  \item \textbf{Soft resampling:} lets some gradient information survive resampling through the weights.
  \item \textbf{Reparametrization trick:} write perturbations and proposals as differentiable functions of Gaussian noise,
  \[
    \delta'(y_t,\lambda)=\Sigma(y_t,\lambda)u+\mu(y_t,\lambda),
    \qquad
    \theta''_j(y_t,\lambda)=\Sigma_t(y_t,\lambda)u_j+\hat\theta_t(y_t,\lambda).
  \]
  \item \textbf{Scibior-Wood correction:} replace the new weights by the surrogate
  \[
    \tilde w'_j \leftarrow w'_j\,\frac{q_{\phi(j)}}{\mathrm{sg}[q_{\phi(j)}]},
  \]
  which leaves the forward pass unchanged but injects missing backward terms.
\end{enumerate}
\vspace{0.4em}
For the MSE loss, the paper shows that this correction reproduces the needed score-function terms.
\end{frame}

\begin{frame}{Appendix D.1--D.3: Bayesian loss and policy-gradient correction}
For continuous estimation tasks, the single-run loss is an MSE-type quantity,
\[
  \ell(\hat\theta,\theta)=\mathrm{tr}\!\left[G(\hat\theta-\theta)(\hat\theta-\theta)^\top\right],
\]
while discrete tasks use a $0$-$1$ loss.
\vspace{0.4em}
Monte Carlo over a batch of $B$ simulations gives
\[
  L(\lambda)=\frac{1}{B}\sum_{k=1}^B \ell(\hat\theta_k,\theta_k).
\]
Because outcomes are sampled from a policy-dependent distribution, the correct training loss becomes
\[
  \tilde\ell(\omega_k,\lambda)=\ell(\omega_k,\lambda)
  +\mathrm{sg}[\ell(\omega_k,\lambda)]\log p(y_{k,0:M-1}\mid\theta_k,\lambda),
\]
so that the gradient includes both pathwise terms and score-function terms.
\end{frame}

\begin{frame}{Appendix D.4--D.6: cumulative and logarithmic losses}
A final-time loss alone may not yield good performance at intermediate resource budgets.
\vspace{0.5em}
\begin{block}{Cumulative loss}
\[
  L_{\mathrm{cum}}(\lambda)=\frac{1}{MB}\sum_{t=0}^{M-1}\sum_{k=1}^B
  \ell(\hat\theta_{k,t},\theta_k).
\]
This pressures the policy to use resources sensibly \emph{throughout} the experiment.
\end{block}
\begin{block}{Logarithmic loss}
Taking a logarithm of the batch-averaged loss stabilizes optimization when the precision spans several orders of magnitude.
\end{block}
\begin{itemize}
  \item A baseline can be subtracted inside the stop-gradient term to reduce gradient variance.
  \item The paper uses Adam rather than plain SGD for the actual parameter updates.
\end{itemize}
\end{frame}

\begin{frame}{Appendix E: simulation practice, hyperparameters, and scaling}
\textbf{Practical tuning strategy}
\begin{enumerate}
  \item Choose the number of particles $N$ large enough that posterior discretization is not the bottleneck.
  \item Use the remaining memory budget to set the batch size $B$.
  \item Tune the initial learning rate with the chosen loss (cumulative or logarithmic).
\end{enumerate}
\vspace{0.5em}
\textbf{Complexity trends reported in Appendix E.3}
\[
  \text{Memory}\sim \mathcal O(dN+N+n_\ell n_u),
  \qquad
  \text{Time}_{\mathrm{seq}}\sim \mathcal O(MN+Mn_\ell n_u^2),
\]
for parameter dimension $d$, measurements $M$, and a neural network with $n_\ell$ layers and $n_u$ units per layer.
\begin{itemize}
  \item Training benefits strongly from GPU parallelism over the batch.
  \item Deployment can be far cheaper, especially for non-adaptive strategies or offline Bayesian updates.
\end{itemize}
\end{frame}

\begin{frame}{Applications highlighted in the main text}
\textbf{NV-center DC magnetometry}
\begin{itemize}
  \item Control variable: Ramsey free-evolution time $\tau$.
  \item Likelihood: binary outcomes with dephasing-dependent contrast.
  \item Learned strategies beat common heuristics and prior model-free RL baselines in the reported settings.
\end{itemize}
\vspace{0.5em}
\textbf{Agnostic Dolinar receiver}
\begin{itemize}
  \item Task: discriminate coherent states while also treating the amplitude as an unknown nuisance parameter.
  \item Control variable: beam-splitter reflectivity/phase sequence.
  \item Adaptive RL approaches the finite-reference Helstrom limit more closely than prior hand-designed strategies.
\end{itemize}
\vspace{0.4em}
\textit{Figures intentionally omitted here; the corresponding plots from the article can be inserted later.}
\end{frame}

\begin{frame}{Appendix F: frequentist optimization via Fisher information}
The library also supports local frequentist design by optimizing the Cramer-Rao bound.
\vspace{0.4em}
\[
  K=\mathbb E\!\left[(\hat\theta-\theta)(\hat\theta-\theta)^\top\right]\ge F^{-1}(\theta),
\]
with trajectory-level Fisher information
\[
  F_{ij}(\theta)=\mathbb E_y\!\left[\frac{\partial\log p(y\mid x,\theta)}{\partial\theta_i}
  \frac{\partial\log p(y\mid x,\theta)}{\partial\theta_j}\right].
\]
Contracting with $G\ge 0$ gives the scalar loss
\[
  L(\lambda)=\mathrm{tr}\!\left(GF^{-1}\right).
\]
\begin{itemize}
  \item The gradient requires differentiating the Fisher matrix through the policy.
  \item Importance sampling can be used if rare trajectories dominate the Fisher information.
\end{itemize}
\end{frame}

\begin{frame}{Appendix G: lower bounds and how to benchmark learned policies}
Appendix G builds Bayesian Cramer-Rao style lower bounds tailored to the examples.
\vspace{0.4em}
\begin{itemize}
  \item For measurement-limited protocols,
  \[
    \Delta^2\hat\theta\ge \frac{1}{\sup_h \mathbb E_\theta[I(\theta\mid \tau_h)] + I(\pi)}
    \ge \frac{1}{M\,\mathbb E_\theta[\sup_\tau I(\theta\mid\tau)] + I(\pi)}.
  \]
  \item For time-limited protocols, the relevant object is the \emph{normalized} Fisher information $I(\theta\mid\tau)/\tau$.
  \item In NV magnetometry, binary outcomes also imply bit-counting limits that can dominate when the Fisher-information bound is loose.
\end{itemize}
\vspace{0.4em}
These bounds define the gray ``impossible region'' used in the paper's benchmark plots.
\end{frame}

\begin{frame}{Appendix H: backward recursion and the control-theory viewpoint}
Appendix H reframes the learned policy as an approximation to a dynamic-programming solution.
\vspace{0.4em}
\begin{itemize}
  \item At the last step, the optimal control $x^*_{M-1}$ minimizes the expected terminal loss conditioned on the current PF ensemble.
  \item Recursing backward defines a family of optimal decision rules $r_t(p_t)$ acting on the particle-filter state.
  \item The optimization is \textbf{non-myopic}: a control at time $t$ is chosen for its effect on the \emph{final} loss, not only the next measurement.
\end{itemize}
\vspace{0.5em}
\begin{block}{Interpretation}
The neural network is not a black box replacing physics; it is a function approximator for the intractable backward-recursion solution on belief states.
\end{block}
\end{frame}

\begin{frame}{Appendix I: backpropagation structure and the stop-gradient trick}
The terminal loss can be written as a composition of Bayes-update maps acting on the initial PF ensemble.
\vspace{0.5em}
The derivative expands into repeated products of two kinds of terms:
\begin{itemize}
  \item direct dependence of the Bayes update on the control selected by the agent,
  \item indirect dependence through earlier PF states fed back into later policy evaluations.
\end{itemize}
\vspace{0.5em}
Appendix I shows that:
\begin{itemize}
  \item these recursive terms generate expensive higher-order gradient paths,
  \item inserting a stop-gradient on the PF input,
  \[
    x_{t+1}=F_\lambda\big(\mathrm{sg}[P(\theta\mid x_t,y_t);R_t;t]\big),
  \]
  removes a large fraction of backward complexity,
  \item yet the resulting optimization remains non-myopic because each local update still targets final precision.
\end{itemize}
\end{frame}

\begin{frame}{Take-home messages}
\begin{enumerate}
  \item \textbf{Belief-state control:} quantum metrology can be formulated as sequential control on a Bayesian belief state.
  \item \textbf{Differentiable inference:} particle filtering can be embedded into model-aware RL through soft resampling, reparametrization, and score-function corrections.
  \item \textbf{Versatility:} the same framework supports adaptive/non-adaptive and Bayesian/frequentist design.
  \item \textbf{Interpretability:} Appendices H and I reveal the learned policy as an approximation to backward recursion with tractable backpropagation.
\end{enumerate}
\vspace{0.6em}
\textbf{Possible closing slide additions:}
benchmark plots, platform-specific diagrams, and an implementation workflow from \texttt{qsensoropt}.
\end{frame}



\end{document}
